\element{diffusionPart}{
  \begin{question}{equationDiffusionPart}\bareme{1}
    L'équation de la diffusion de particules s'écrit
    \begin{multicols}{2}
      \begin{reponses}
        \bonne{$\frac{\partial n}{\partial t}-D_{th}\Delta n=0$}
        \mauvaise{$\frac{\partial^2 n}{\partial t^2}+D_{th}\Delta n=0$}
        \mauvaise{$\frac{\partial^2 n}{\partial t^2}-D_{th}\Delta n=0$}
        \mauvaise{$\frac{\partial^2 n}{\partial t^2}-D_{th}\Delta^2 n=0$}
      \end{reponses}
    \end{multicols}
  \end{question}
}

\element{reversibiliteDiffusion}{
  \begin{question}{reversibiliteEquationDiffusion}\bareme{1}
    L'équation de diffusion est réversible.
    \begin{multicols}{2}
      \begin{reponses}
        \mauvaise{Vrai}
        \bonne{Faux}
      \end{reponses}
    \end{multicols}
  \end{question}
}

\element{fick}{
  \begin{questionmult}{loiFick}\bareme{haut=2,p=0}
    Dans la loi de Fick,
      \begin{reponses}
        \bonne{Le signe "-" indique que les particules vont de là où il y en a le plus vers là où il y en a le moins.}
        \bonne{Le terme $D$ se mesure en \SI{}{m^{2}s^{-1}}}
        \bonne{Le vecteur densité de courant particulaire se mesure en \SI{}{m^{-2}s^{-1}}}
      \end{reponses}
  \end{questionmult}
}
